\documentclass[11pt,a4paper]{../_shared/altacv}

\geometry{left=1.5cm,right=1.5cm,top=1.cm,bottom=1.2cm,columnsep=1.5cm}

\usepackage[utf8]{inputenc}
\usepackage[T1]{fontenc}
\usepackage[french]{babel}
\usepackage{paracol}
\usepackage{xcolor}
\usepackage{hyperref}
\usepackage{fontawesome5}
\usepackage{setspace}

\definecolor{SlateGrey}{HTML}{2E2E2E}
\definecolor{LightGrey}{HTML}{666666}
\definecolor{DarkPastelRed}{HTML}{8AC0D9}
\definecolor{PastelRed}{HTML}{00B0DA}
\definecolor{GoldenEarth}{HTML}{B4AD0C}
\colorlet{name}{black}
\colorlet{tagline}{PastelRed}
\colorlet{heading}{DarkPastelRed}
\colorlet{headingrule}{GoldenEarth}
\colorlet{subheading}{PastelRed}
\colorlet{accent}{PastelRed}
\colorlet{emphasis}{SlateGrey}
\colorlet{body}{LightGrey}

\renewcommand{\familydefault}{\sfdefault}

\name{\Large Guillaume BOILEAU}
\tagline{AI \& DSP Engineer | Signal Processing, Noise Modeling, Machine Learning \& Algorithm Validation}
\personalinfo{
  \email{guillaume.boileaupro@gmail.com}
  \phone{+33 6 59 94 85 84}
  \location{Alpes-Maritimes, France}
  \linkedin{guillaume-boileau-phd-891b81265}
  \researchgate{Guillaume-Boileau-2}
  \orcid{0000-0002-3576-6968}
}

\setlength{\parskip}{0.75\baselineskip}

\begin{document}
\makecvheader
\vspace{-0.6cm}

\cvsection{Cover Letter}
\vspace{-0.4cm}

{\color{accent}\textbf{Subject:}} \textit{Application for the AI\&DSP Engineer position -- Pulse Audition / EssilorLuxottica}

{\color{body}

Dear Madam, Dear Sir,

I am writing to apply for the AI\&DSP Engineer position within Pulse Audition at EssilorLuxottica in Sophia Antipolis. This opportunity immediately caught my attention because it combines several dimensions that are central to my background and to the direction I would like to give to my career: advanced signal processing, statistical modeling, machine learning, algorithm prototyping, performance validation and product-oriented R\&D applied to a concrete human need.

The mission of Pulse Audition is particularly meaningful to me. Developing AI and DSP technologies for smart eyewear and assistive hearing applications is not only a technically ambitious challenge, but also a project with direct social impact. Improving communication, autonomy and inclusion for people affected by hearing loss is the kind of applied innovation that gives strong purpose to research and engineering work. I am therefore highly motivated by the possibility of contributing to a team that brings together AI, audio signal processing, audiology, hardware and software expertise to create technologies that can be integrated into real products.

My background has been built around the analysis, modeling and extraction of weak signals embedded in complex noisy environments. During my PhD and subsequent R\&D activities, I worked extensively on gravitational-wave data analysis, where the central challenge is to identify, separate and characterize extremely low signal-to-noise signals in the presence of multiple noise sources, overlapping components and strong uncertainty. This experience gave me a solid foundation in digital signal processing, spectral analysis, statistical inference, Bayesian modeling, simulation workflows and algorithm validation.

Although my previous applications were not audio-specific, the underlying technical challenges are strongly connected to those encountered in audio and speech processing. In both cases, the objective is to extract useful information from noisy and complex data, separate relevant signals from unwanted components, model the spectral structure of the signal and noise, assess the robustness of algorithms, and validate performance under realistic conditions. This is why I see a strong continuity between my previous work on low-SNR scientific signals and the development of AI/DSP algorithms for speech enhancement, noise reduction and assistive hearing applications.

At the Laboratoire Lagrange, Observatoire de la Côte d'Azur, in connection with CNES activities for the LISA ESA/NASA space mission, I contributed to the development and validation of complex signal-processing and analysis workflows. LISA is a large-scale international space mission involving demanding requirements in terms of modeling, signal analysis, data processing and scientific validation. In this context, I developed CATpyLISA, a Python-based platform dedicated to model integration, comparison and validation of analysis chains. This work required me to design reproducible workflows, evaluate algorithmic outputs, compare models, investigate unexpected behaviours, assess performance limitations and document technical results for collaborative review.

This experience is directly relevant to the AI\&DSP Engineer role because it involved moving from algorithmic concepts to implemented, tested and validated workflows. I was not only working on theoretical aspects of signal modeling, but also on the practical questions that determine whether a method can be trusted and reused: Is the output consistent? Is the result robust? Can the workflow be reproduced? Can the source of an anomaly be identified? Can different contributors understand and use the tool? These questions are also essential when developing AI/DSP algorithms intended to be integrated into embedded or wearable devices.

Previously, at the University of Antwerp, I worked on noise modeling and performance studies for precision instruments. I developed analysis pipelines to identify, characterize and mitigate environmental and instrumental noise sources affecting measurement quality and system sensitivity. This work strengthened my ability to interpret complex measurement data, understand how noise sources degrade performance, and build robust analysis strategies to support system-level diagnostics. I believe this background is highly transferable to audio applications, especially for problems such as noise reduction, signal enhancement, robustness evaluation and performance assessment in real-world environments.

My doctoral work at ARTEMIS, Observatoire de la Côte d'Azur, also contributed strongly to my expertise in signal processing and statistical modeling. I worked on the separation of overlapping low-SNR signals, frequency-domain analysis, Bayesian inference and simulation of complex signal environments. These activities required a strong understanding of Fourier-domain methods, probability, statistics, linear algebra and numerical implementation. They also trained me to approach difficult signal-processing problems with rigor, patience and a strong validation mindset.

From a technical point of view, I have strong experience in Python for algorithm development, data analysis, numerical modeling, simulation and validation workflows. I have also used C in my academic and technical background, and I am comfortable working in Linux environments with Git/GitLab, Docker, scripting and structured software workflows. I am familiar with PyTorch and machine learning concepts, and I am particularly motivated to deepen their application to audio-related models, hybrid AI/DSP pipelines and embedded deployment constraints.

I am aware that the AI\&DSP Engineer role requires specific audio expertise, including speech enhancement, feedback suppression, beamforming, acoustic systems and embedded audio constraints. My background is not that of a pure audio engineer, and I do not want to overstate this point. However, I bring a strong foundation in the core disciplines required to grow efficiently in this direction: digital signal processing, noise modeling, statistical inference, algorithm evaluation, Python prototyping, performance validation and scientific literature analysis. I am used to entering technically demanding domains, understanding their constraints quickly, and building reliable methods from complex requirements.

I also bring recent industrial software experience from VEV Platform Services France, where I worked on software services connected to operational hardware systems in the electric mobility sector. This experience helped me move closer to product and operational constraints. I worked on data exchanges, service reliability, hardware/software interactions, issue analysis and production-oriented technical follow-up. This complements my research background by giving me a more concrete view of software robustness, maintainability and operational impact.

What I find especially attractive in your team is the balance between exploration and product integration. The role is not limited to research; it also requires adapting, optimizing and productizing algorithms for embedded deployment. This is exactly the kind of transition I would like to contribute to: transforming advanced signal-processing and AI methods into robust, testable and usable technologies. I am particularly interested in the challenges of real-time processing, on-device constraints, model adaptation and the interaction between algorithm design, hardware limitations and user experience.

Beyond the technical match, I believe I would fit well into a multidisciplinary team such as Pulse Audition. Throughout my career, I have worked in international collaborations involving researchers, engineers, software developers and technical contributors from different backgrounds. I have coordinated technical tasks, supervised interns, documented complex results and contributed to collective validation efforts. I value clear communication, reproducibility, structured work and constructive collaboration, especially in environments where different areas of expertise must converge toward a common product goal.

Joining EssilorLuxottica would represent for me a strong opportunity to apply my experience in signal processing and algorithm validation to a field with direct industrial and societal impact. I am particularly motivated by the idea of contributing to smart eyewear technologies that combine AI, DSP, embedded systems and human-centered innovation. I would bring a rigorous scientific mindset, strong analytical skills, practical software experience and a genuine motivation to grow into audio AI and embedded DSP applications.

I would welcome the opportunity to discuss how my background in low-SNR signal analysis, noise modeling, statistical inference, Python prototyping and validation of complex processing chains could contribute to the development of AI\&DSP algorithms for Pulse Audition.

Yours faithfully,

\textbf{Guillaume Boileau}

}

\end{document}